%% Appendix: stain-normalization granularity (MedIA revision, v2 pipeline).
%% Final numbers: results/v2/analysis/{C,S}/{cohort}/summary_id.csv (snapshot 1381bdd759f9),
%% ID-validation selection. Pre-fix per-tile numbers: $V2/runs_quarantine_macenko_prefix (tier C,
%% c17 fold 4). Tile-makeup diagnostic: revision/macenko_resolution.md section 4 ($V2/runs_diag/macgroup).
\section{Stain-normalization granularity}
\label{app:macenko}

\paragraph{Per-slide and per-tile estimation} Our Macenko arms estimate one
stain matrix per slide and scanner (Section~\ref{M-sec:setup-arms}). The per-tile
variant runs the same estimator on every tile separately, with the
non-negativity clamp, and maps it to a reference fitted on 256 training tiles.
A tile with less than 5\% tissue keeps the reference matrix. Both variants use
no labels. Both were run with the full protocol: 5 seeds at the compact tier,
3 at the standard tier, all 20 folds.

\paragraph{Result} Table~\ref{tab:macenko} gives the worst-domain AUROC of each
variant. The two differ mainly on Camelyon17 center~4. There, per-tile
normalization gives an off-site AUROC of 0.524 at the compact tier and 0.557 at
the standard tier. With per-slide normalization the scores are 0.952 and 0.976.
The model is not the cause: standard-tier ERM without normalization scores 0.977
on center~4. On the other four Camelyon17 centers, the per-tile variant is
within 0.018 of the per-slide one at the compact tier and equal or better at
the standard tier, by up to 0.034. On canine cSCC, per-slide normalization is
higher by 0.027 (compact) and 0.036 (standard). On the two MIDOG cohorts the
variants differ by at most 0.026. On Camelyon17 the choice of granularity alone moves the arm from rank 4 of 20
to last at the compact tier, and from rank 10 of 21 to last at the standard
tier. The failing domain has the palest tissue: center-4 tiles are 27--34\% tissue pixels,
against 58--83\% in the other centers. A per-tile estimate rests on the
1,100--1,400 tissue pixels of a 64~px tile, usually of a single tissue type.

\begin{table}[width=\linewidth,pos=t]
\caption{Worst-domain AUROC (ID-validation selection) of the stain-normalization
arms, with ERM and H\&E jitter for reference, and the off-site AUROC on
Camelyon17 center~4 (C17 c4). Means over 5 seeds (compact tier) and 3 seeds
(standard tier). $^{\dagger}$Transductive.}
\label{tab:macenko}
\footnotesize
\setlength{\tabcolsep}{3pt}
\begin{tabular*}{\tblwidth}{@{\extracolsep{\fill}}lcccccccccc@{}}
\toprule
 & \multicolumn{5}{c}{Compact tier} & \multicolumn{5}{c}{Standard tier} \\
\cmidrule(lr){2-6}\cmidrule(lr){7-11}
Arm & C17 c4 & C17 & Canine & MIDOG21 & MIDOG++ & C17 c4 & C17 & Canine & MIDOG21 & MIDOG++ \\
\midrule
ERM & 0.698 & 0.698 & 0.551 & 0.765 & 0.766 & 0.977 & 0.895 & 0.534 & 0.847 & 0.870 \\
H\&E jitter & 0.917 & 0.904 & 0.674 & 0.864 & 0.809 & 0.985 & 0.958 & 0.850 & 0.875 & 0.880 \\
Macenko (slide)$^{\dagger}$ & 0.952 & 0.904 & 0.698 & 0.870 & 0.828 & 0.976 & 0.915 & 0.819 & 0.860 & 0.873 \\
TIA-style$^{\dagger}$ & 0.982 & 0.889 & 0.730 & 0.864 & 0.810 & 0.960 & 0.934 & 0.823 & 0.871 & 0.877 \\
Macenko (tile) & 0.524 & 0.524 & 0.671 & 0.844 & 0.816 & 0.557 & 0.557 & 0.783 & 0.846 & 0.877 \\
\bottomrule
\end{tabular*}
\end{table}

\paragraph{Negative concentrations} The earlier version of this pipeline used
per-tile estimation with unconstrained least-squares concentrations, as in
Macenko's reference code. In a 64~px tile the estimated hematoxylin and eosin
vectors are often nearly collinear (cosine 0.95--0.98). Eosin-rich pixels then
receive large negative hematoxylin values and come out saturated pink after
reconstruction. On center~4 this artifact affected 90\% of tumor tiles but
54\% of normal tiles. The off-site AUROC on center~4 was 0.42 (5 seeds, compact
tier), below every arm of Table~\ref{tab:macenko}. Clamping the concentrations at zero removes the
artifact but not the unreliable per-tile stain matrix. The per-slide estimate
removes both. The earlier records are kept apart and are not used in any table.

\paragraph{Dependence on the tile set} The per-slide estimate of a test slide
uses that slide's own unlabeled tiles. In Camelyon17 those tiles are the
benchmark's patches, sampled from annotated regions and class-balanced. The
share of tumor tiles per test slide therefore ranges from 0 to 0.93. To bound
the effect of this makeup, we fixed the trained weights of five compact-tier
runs on center~4 (seeds 0--4, final configuration) and re-estimated each test
slide's stain matrix from different subsets of its tiles. Labels were used only
to form these diagnostic subsets. The mean off-site AUROC was 0.945 with all
tiles, 0.948 with only negative tiles and 0.934 with only positive tiles: a
spread of 0.014. With no per-slide estimate, using only the reference, it fell
to 0.560. For TIA-style the corresponding values were 0.981, 0.972, 0.978 and
0.897. The trained models therefore rely on the per-slide normalization, but
not much on which tiles of the slide feed it. A label-independent tissue sample
of each whole-slide image would match deployment more closely. It is not
available, because the cohorts are distributed as patches.
